\documentclass[../../../../main.tex]{subfiles}

\begin{document}

\section*{東大理科数学 1980年}

\begin{itemize}
    \item 所要時間: 80.5 分 / 120 分
\end{itemize}

\begin{center}
\begin{tabular}{|c|c|c|c|c|}
\hline
問 & 分野 & 計 & 見 & 統 \\ \hline
1 & 図形 & C & B & B \\ \hline
2 & 空間 & B & B & B \\ \hline
3 & 行列 & & & \\ \hline
4 & 多変数 & C & B & B \\ \hline
5 & 場合の数 & B & A & A \\ \hline
6 & 多変数 & C & B & B \\ \hline
\end{tabular}
\end{center}

\begin{itemize}
    \item 1で検算中にミス
    \item 3で, $(\pi/2 \text{ と } \pi/4)$, $(\pi/3 \text{ と } 2\pi/3)$ をまちがえる $\Rightarrow$ よくやるので注意
    \item 4の因数分解ミスはゆるされない $\rightarrow$ 細心の注意＋おいた変数とは別のほうにかくな
\end{itemize}

\newpage

\subsection*{第1問}

\noindent
\textbf{[解]} $0 < x < \frac{1}{2} \quad \cdots \text{①}$

\begin{center}
\begin{tikzpicture}[scale=2.0]
  % Triangle ABC
  \coordinate (A) at (60:3);
  \coordinate (B) at (0,0);
  \coordinate (C) at (3,0);

  % Points P, Q, R on sides
  \coordinate (P) at (0.8,0);
  \coordinate (Q) at ($(C)!0.267!(A)$);
  \coordinate (R) at ($(A)!0.267!(B)$);

  % Intersections
  \coordinate (A') at (intersection of A--P and C--R);
  \coordinate (B') at (intersection of A--P and B--Q);
  \coordinate (C') at (intersection of B--Q and C--R);

  % Draw main triangle
  \draw[thick] (A) -- (B) -- (C) -- cycle;

  % Draw lines AP, BQ, CR
  \draw (A) -- (P);
  \draw (B) -- (Q);
  \draw (C) -- (R);

  % Node labels
  \node[above] at (A) {$A$};
  \node[below left] at (B) {$B$};
  \node[below right] at (C) {$C$};

  \node[below] at (P) {$P$};
  \node[right] at (Q) {$Q$};
  \node[left] at (R) {$R$};

  \node[above right] at (A') {$A'$};
  \node[left] at (B') {$B'$};
  \node[below] at (C') {$C'$};

  % Length labels
  \node[below] at ($(B)!0.5!(P)$) {$x$};
  \node[right] at ($(C)!0.5!(Q)$) {$x$};
  \node[left] at ($(A)!0.5!(R)$) {$x$};

  \node[right] at ($(A)!0.5!(A')$) {$p$};
  \node[right] at ($(A')!0.5!(B')$) {$r$};
  \node[right] at ($(B')!0.5!(P)$) {$q$};

  \node[above] at ($(B)!0.5!(B')$) {$p$};
  \node[above] at ($(B')!0.5!(C')$) {$r$};
  \node[above] at ($(C')!0.5!(Q)$) {$q$};

  \node[left] at ($(C)!0.5!(C')$) {$p$};
  \node[left] at ($(C')!0.5!(A')$) {$r$};
  \node[left] at ($(A')!0.5!(R)$) {$q$};
\end{tikzpicture}
\end{center}

\begin{enumerate}
  \item[(1)] 対称性から $\overline{BB'} = \overline{CC'} = \overline{AA'} = p$,
  $\overline{PB'} = \overline{QC'} = \overline{RA'} = q$, $\overline{AB'} = \overline{BC'} = \overline{CA'} = r$
  とおける。メネラウスの定理から
  \[
  \frac{B'C'}{BB'} \cdot \frac{A'R}{C'A'} \cdot \frac{AB}{RA} = 1 \quad \cdots \text{②}
  \]
  \[
  \frac{QA}{CQ} \cdot \frac{B'P}{AB'} \cdot \frac{BC}{PB} = 1 \quad \cdots \text{③}
  \]
  まず ③ から
  \[
  \frac{1-x}{x} \cdot \frac{q}{p+r} \cdot \frac{1}{x} = 1
  \]
  ② から
  \[
  \frac{r}{p} \cdot \frac{q}{r} \cdot \frac{1}{x} = 1
  \]
  したがって
  \[
  \begin{cases}
  q = xp \\
  r = \frac{1-2x}{x} p
  \end{cases}
  \quad \cdots \text{④}
  \]
  一方, $\triangle BCR$ に余弦定理を用いて,
  \[
  p + q + r = \sqrt{x^2 - x + 1} \quad \cdots \text{⑤}
  \]
  ④ を ⑤ に代入して,
  \[
  \left( 1 + x + \frac{1-2x}{x} \right) p = \sqrt{x^2 - x + 1} \quad \therefore p = \frac{x}{\sqrt{x^2 - x + 1}}
  \]
  ④ から
  \[
  r = \frac{1-2x}{\sqrt{x^2 - x + 1}}, \quad q = \frac{x^2}{\sqrt{x^2 - x + 1}}
  \]
  以上から, $\overline{BB'} = p = \frac{x}{\sqrt{x^2 - x + 1}}$, $\overline{PB'} = q = \frac{x^2}{\sqrt{x^2 - x + 1}}$

  \item[(2)] $\triangle A'B'C'$ は一辺の長さ $r$ の正三角形でその面積 $T$ は
  \[
  T = \frac{\sqrt{3}}{4} r^2 = \frac{\sqrt{3}}{4} \cdot \frac{4x^2 - 4x + 1}{x^2 - x + 1}
  \]
  と表せる。これが $\frac{1}{2} \cdot \frac{\sqrt{3}}{4}$ に等しいので
  \[
  \frac{1}{2} = \frac{4x^2 - 4x + 1}{x^2 - x + 1}
  \]
  \[
  8x^2 - 8x + 2 = x^2 - x + 1
  \]
  \[
  7x^2 - 7x + 1 = 0
  \]
  \[
  \therefore x = \frac{7 \pm \sqrt{21}}{14}
  \]
  $0 < x < \frac{1}{2}$ から複号負を採用して,
  \[
  x = \frac{7 - \sqrt{21}}{14}
  \]
\end{enumerate}

\subsection*{第2問}

\noindent
\textbf{[解]} $NA, NB$ と $k$ の交点を $C, D$ とする。$NP$ と $k$ の交点を $Q$ とする。

$AB$ 上を $P$ がうごく時, $Q$ は $C, D$ を両端とする大円の劣弧上を動く。この長さは, 半径 1 の四分円の弧長にひとしく,
\[
\frac{1}{4} \cdot 2\pi = \frac{1}{2}\pi \quad \cdots \text{①}
\]

次に $AB$ 上を $P$ がうごく時, 平面 $ABN$ で $k$ を切ると下図で $\text{+++}$ 上を $Q$ がくまなくうごく。この円は, 1 辺の長さ $2\sqrt{2}$ の正三角形の外接円で, その半径 $\frac{2\sqrt{6}}{3}$ だから
\[
\stackrel{\frown}{CD} = \frac{1}{3} \cdot \frac{2\sqrt{6}}{3} \pi = \frac{2\sqrt{6}}{9}\pi \quad \cdots \text{②}
\]

\begin{center}
\begin{tikzpicture}[scale=1.3]
  % Diagram 1: Sphere sketch
  \begin{scope}[shift={(0,0)}]
    \draw (0,0) circle (1.2);
    \draw (-1.2,0) arc (180:360:1.2 and 0.4);
    \draw[dashed] (-1.2,0) arc (180:0:1.2 and 0.4);
    
    \coordinate (N) at (0,1.2);
    \coordinate (A) at (-1.8,-1.5);
    \coordinate (B) at (1.8,-1.5);
    
    \draw (N) -- (A);
    \draw (N) -- (B);
    \draw (A) -- (B);
    
    \node[above] at (N) {$N$};
    \node[below left] at (A) {$A$};
    \node[below right] at (B) {$B$};
  \end{scope}

  % Diagram 2: Triangle ABN section
  \begin{scope}[shift={(5,0)}]
    \coordinate (N) at (0, 2.45);
    \coordinate (A) at (-2, 0);
    \coordinate (B) at (2, 0);

    \draw[thick] (N) -- (A) -- (B) -- cycle;

    \coordinate (C) at ($(N)!0.5!(A)$);
    \coordinate (D) at ($(N)!0.5!(B)$);

    % Outer circle / arc CD
    \draw[thick] (C) arc (210:330:2.3);

    \node[above] at (N) {$N$};
    \node[below left] at (A) {$A$};
    \node[below right] at (B) {$B$};
    \node[left] at (C) {$C$};
    \node[right] at (D) {$D$};

    \node[left] at ($(N)!0.5!(C)$) {$\sqrt{2}$};
    \node[left] at ($(C)!0.5!(A)$) {$\sqrt{2}$};
    \node[right] at ($(N)!0.5!(D)$) {$\sqrt{2}$};
    \node[right] at ($(D)!0.5!(B)$) {$2\sqrt{2}$};
    \node[below] at ($(A)!0.5!(B)$) {$2\sqrt{2}$};

    \node[below] at (0,1.0) {$\frac{2}{3}\pi$};
  \end{scope}
\end{tikzpicture}
\end{center}

以上の ①, ② から, もとめる長さは, (これらが $C, D$ 以外で交わらないことから)
\[
\frac{1}{2}\pi + \frac{2\sqrt{6}}{9}\pi = \left(\frac{1}{2} + \frac{2\sqrt{6}}{9}\right)\pi
\]

\subsection*{第3問}

(未解答)

\subsection*{第4問}

\noindent
\textbf{[解]} $S = \sin t, C = \cos t$ とする。 $\overline{OP}^2 = f(t)$ とする。
\[
\begin{aligned}
f(t) &= (S+C)^2 + k^2 (S^2 C)^2 \\
&= (S+C)^2 + k^2 S^4 C^4 \\
&= k^2 p^4 + 2p + 1 \quad \left( p = SC \text{ とすると } -\frac{1}{2} \le p \le \frac{1}{2} \right)
\end{aligned}
\]
\[
\frac{df}{dp} = 4k^2 p^3 + 2 = 4k^2 \left(p^3 + \frac{1}{2k^2}\right) = 4k^2 (p+\alpha)(p^2 - p\alpha + \alpha^2) \quad \left( \alpha = \sqrt[3]{\frac{1}{2k^2}} \right)
\]
から, 下表を得る。 ($\because k > 0$)

\begin{enumerate}
  \item[$1^\circ$] $\frac{1}{2} \ge \alpha$ の時 ($2 \le k$)

  \begin{center}
  \begin{array}{c|c|c|c|c|c}
  p & -\frac{1}{2} & \dots & -\alpha & \dots & \frac{1}{2} \\ \hline
  f' & & - & 0 & + & \\ \hline
  f & & \searrow & \text{極小} & \nearrow &
  \end{array}
  \end{center}

  \[
  \begin{cases}
  f\left(\frac{1}{2}\right) = \frac{k^2}{16} + 2 \\
  f\left(-\frac{1}{2}\right) = \frac{k^2}{16} \\
  f(-\alpha) = -\frac{3}{2}\alpha + 1
  \end{cases}
  \]

  表及び $f$ の値から
  \[
  \begin{cases}
  \max f = f\left(\frac{1}{2}\right) = \frac{1}{16}k^2 + 2 \\
  \min f = f(-\alpha) = 1 - \frac{3}{2} \sqrt[3]{\frac{1}{2k^2}}
  \end{cases}
  \]

  \item[$2^\circ$] $\frac{1}{2} \le \alpha$ の時 ($k \le 2$)

  $f'(p) \ge 0$ より, $f$ は $p$ について単調増加だから,
  \[
  \begin{cases}
  \max f = f\left(\frac{1}{2}\right) = \frac{1}{16}k^2 + 2 \\
  \min f = f\left(-\frac{1}{2}\right) = \frac{1}{16}k^2
  \end{cases}
  \]
\end{enumerate}

したがって
\[
\begin{cases}
0 < k \le 2 \text{ の時}, & \max f = \frac{1}{16}k^2 + 2, \quad \min f = \frac{1}{16}k^2 \\[6pt]
2 \le k \text{ の時}, & \max f = \frac{1}{16}k^2 + 2, \quad \min f = 1 - \frac{3}{2}\sqrt[3]{\frac{1}{2k^2}}
\end{cases}
\]

\end{document}
